\subsection{The stationary point of the child}\label{sec:B4T02}
Retain the notation of Section~\ref{sec:B4T01}. Now $B$ is the child's RC line--square matrix and $D$ is its square--diagonal incidence matrix. Section~\ref{sec:B3T07} constructs
\[
 \bar p=w+z,\quad B\bar p=\mathbf1,\quad
 a=z/w\in\operatorname{im}B^T,\quad \|a\|_\infty=O(k^{-1}).
\]
Let $d$ be the vector of losses on the retained diagonals and set $u=d-D^Tz$.
The vacancies of $\bar p$ are $h+u\ge\eta/2$.

\paragraph{Stationarity residual after rebalancing.}
By stationarity of the parent and the exact expressions for the logarithms,
\[
 \nabla H_{\rm child}(\bar p)\equiv
 D\log(1+u/h)+a-\log(1+a)-3\log(1-z/(1-w))
 \pmod{\operatorname{im}B^T}.
\]
In particular, the tangent gradient has a representative $e$ such that
\[
 e=D(u/h)+O_\infty(k^{-2}),\qquad \|e\|_\infty=O(k^{-1}).
\]
Correcting the RC masses does not in general cancel the first diagonal term. We do not assert that $p^*-\bar p=O_{\rm rel}(k^{-2})$.

\paragraph{Newton step and uniform bounds.}
Let $Z$ be an orthonormal basis of $\ker B$ and set
\[
 K=\operatorname{diag}(1/\bar p-3/(1-\bar p))
       +D\operatorname{diag}((h+u)^{-1})D^T,\quad
 R=Z(Z^TKZ)^{-1}Z^T,\quad s=Re.
\]
Apply the matrix estimate of Section~\ref{sec:B3T05} to the feasible point $\bar p$ with the inherited gap. The number of rows in the child satisfies $k-1\ge\delta_0 n/2$ for $k\ge2$, so we may use $\delta_0/2$ even at the endpoint $k=\delta_0n$. This gives $\|\operatorname{diag}(\bar p)^{-1}R\|_{\infty\to\infty}\le M$.
Write $\|v\|_{\bar p}=\|v/\bar p\|_\infty$ and $\sigma=\|s\|_{\bar p}=O(k^{-1})$.

For $v\in\ker B$ with $\|v\|_{\bar p}\le r$ sufficiently small,
\[
 \nabla H(\bar p+v)=\nabla H(\bar p)-Kv+\mathcal N(v),
\]
\[
 \|\mathcal N(v)\|_\infty\le Cr^2,\quad
 \|\mathcal N(v)-\mathcal N(w)\|_\infty\le Cr\|v-w\|_{\bar p}.
\]
The constants do not depend on a lower bound for $\bar p$. Indeed, $|(D^Tv)_d|\le r\sum_{e\in d}\bar p_e\le r$, and the expansion of each square logarithm uses the relative increment $v_e/\bar p_e$. An expansion in $v_e$ alone would not give this control.
Stationarity is equivalent to $v=s+R\mathcal N(v)$.
For large $k$, the right-hand side defines a map that takes the ball of radius $2\sigma$ into itself and has Lipschitz constant less than one. If $\sigma=0$, then $\bar p$ is already stationary.
There is therefore a point $p^*$ with
\[
 \boxed{\|(p^*-\bar p)/\bar p\|_\infty=O(k^{-1}),\qquad
 \|(p^*-\bar p-s)/\bar p\|_\infty=O(k^{-2}).}
\]
This point remains feasible and positive, with uniformly bounded-below vacancies and $p^*<1/4$.
The negative of the Hessian is positive definite in the convex region $p<1/4$, so the stationary point is unique on that branch. We do not assert uniqueness outside it. The point $\widehat p=\bar p+s$ already approximates the stationary point of the child with relative error $O(k^{-2})$. All bounds are uniform over the children of a parent in the class.

\paragraph{Action increment from $\bar p$.}
At $\bar p$, the linear term $e^Ts=s^TKs$ does not in general vanish. Combining it with the negative quadratic term of the entropy gives
\[
 \boxed{H(p^*)-H(\bar p)=\tfrac12 e^TRe+O(k^{-2}).}
\]
The leading term is $O(k^{-1})$. The cubic remainder is $O(k\|v\|_{\bar p}^3)=O(k^{-2})$ because $\sum\bar p=k-1$ and there are $O(k)$ diagonals. The errors from replacing $v$ by $s$ are also $O(k^{-2})$. The RC normal components pair to zero with every increment.
